\section{组织与研发流程：把课程观点转成执行机制}

当 data、evaluation 与 systems 同时决定实验成败，团队边界也必须围绕同一任务闭环重新组织。本节把技术依赖转成可执行的协作机制：共享版本化任务集、共同审查成本与失败轨迹，并让模型、评估、平台和产品负责人对同一上线门槛负责。

\subsection{跨职能协同的必要性}
根据课程对 ``data-eval-systems'' 的强调，Agent 团队不能按传统 ``模型组/平台组/应用组'' 完全割裂。至少在关键迭代周期内，数据、评估、系统、产品应共用一套优先级与实验台账。

\begin{importantbox}{最小可执行组织单元}
建议建立 ``训练-评估-系统'' 三人小组：一人负责策略实验，一人负责评估设计，一人负责基础设施瓶颈优化。每周以统一任务集复盘，而不是各自汇报局部指标。
\end{importantbox}

\subsection{实验管理与复盘模板}
\begin{table}[H]
\centering
\begin{tabular}{p{2.5cm}p{4.1cm}p{5.1cm}}
\toprule
\textbf{字段} & \textbf{记录内容} & \textbf{目的} \\
\midrule
假设 & 本次改动要解决的具体问题 & 防止 ``试试也行'' 的随机实验 \\
评估口径 & 闭集+开集指标、样本集版本 & 保证实验可比较 \\
成本账单 & 训练/评估/推理分项资源消耗 & 评估 ROI，避免隐性亏损 \\
失败复盘 & 失败类型与恢复路径分析 & 沉淀下一轮数据与规则 \\
\bottomrule
\end{tabular}
\caption{将课程理念落地为团队实验管理模板}
\end{table}

\begin{knowledgebox}{一句话流程准则}
\textit{先定义成功，再开始训练；先定义失败恢复，再扩大权限。}
\end{knowledgebox}

\subsection{本章小结}
课程的价值不仅是技术路线，更是研发方法。将数据、评估、系统协同固化为流程，才能持续获得可复现增益。


